% Cote 119: offprints, reports, courses, CV and letters (typed and printed part): archive cote 119, site page 184 (PDF page 185).
% typescript. Draft compared: none.
% AI-transcribed by Claude Opus 5.5 (claude-opus-5-5, adaptive thinking, effort high), Message Batches API, batch r119_full (msgbatch_01DwZNm5Y7YsKBWyAN4sFs8w), 2026-11-13; not yet checked by eye.
\scanpage{119/184}
\mspage{- 3 -}

trouver.

7.3. \uline{Propriété de moyenne}

Pour $n = 0$ la relation (7.1.1) donne
\[ (7.3.1) \qquad f(z_o) = \frac{1}{2\pi}\int_0^{2\pi} f(z_o + re^{i\theta})\, d\theta \]
autrement dit la valeur de $f$ au point $z_o$ est la moyenne de $f$ sur le cercle
de centre $z_o$ et de rayon r.

\uline{Définition} : On dit que $f$ fonction continue de l'ouvert $U \subset \mathbb{C}$ dans $\mathbb{C}$ pos-%
sède la propriété de moyenne si pour tout $z_o \in U$ et tout $r > 0$ tel que
$\overline{B(z_o,r)} \subset U$ on a la relation (7.3.1).

\uline{Proposition} : toute fonction holomorphe possède la propriété de moyenne.

Ce sont d'ailleurs pas les seules fonctions à posséder cette proprié-%
té. On voit facilement que si une fonction possède la propriété de moyenne, sa
partie réelle et sa partie imaginaire la possèdent aussi. La partie réelle d'une
fonction holomorphe $f$ a donc la propriété de moyenne et n'est pourtant pas
holomorphe sauf si $f$ est constante. En effet, si la partie réelle $P$ de $f$
est holomorphe on a (4.2) : $\dfrac{\partial P}{\partial y} = i\,\dfrac{\partial P}{\partial x}$ ce qui n'est possible que pour $\dfrac{\partial P}{\partial x} = 0$
et $\dfrac{\partial P}{\partial y} = 0$ soit $P$ constante. Le même raisonnement appliqué à la partie ima-%
ginaire elle aussi holomorphe car égale à : $-i(f-P)$ montre qu'elle est aussi
constante d'où notre assertion.
(On consultera la feuille d'exercices n° 10. Exercices 3 et 4).

7.4. \uline{Théorème (principe du maximum)}

Soit $f$ une fonction continue dans un ouvert $U$ de $\mathbb{C}$. Si $f$ possè-%
de la propriété de moyenne et si $|f|$ possède un maximum relatif en un point
$a \in U$ (i.e. si $|f(z)| \leq |f(a)|$ pour tout $z$ dans un voisinage de $a$) alors
$f$ est constante dans un voisinage de $a$.

\uline{Démonstration} :

Le lecteur l'a déjà faite en T.D. (feuille 12 exercice 13). Ne lui

\margin{184}
